\section{Procedencia y reproducción de las series de seguridad}\label{anexo:infosegura}
Las seis series de PNUD/Infosegura (2021) complementan MICS con indicadores
de seguridad y condiciones de vida. Los CSV, el manifiesto de procedencia y
el programa R permiten repetir los cálculos sobre agregados secundarios,
sin registros individuales.

\begin{table}[htbp]
\centering\small
\caption{Series de seguridad: archivos, dimensiones y restricciones}\label{tab:traza_infosegura}
\begin{tabularx}{\textwidth}{@{}p{3.5cm}p{2.5cm}Y@{}}\toprule
Archivo, prefijo \var{infosegura_} & Dimensión & Alcance y restricción\\\midrule
\var{homicidios.csv} & Violencia letal & Tasas por población de referencia; 2021 es proyección.\\
\var{victimizacion.csv} & Experiencia delictiva & Delincuencia en doce meses; universo etario no explicitado.\\
\var{preocupaciones.csv} & Opinión & COVID-19 solo consta en 2021; antes permanece vacío, no cero.\\
\var{denuncias.csv} & Demanda institucional & Promedios mensuales; corte de 2021 no precisado.\\
\var{privaciones.csv} & Condiciones de vida & Cambios 2019--2020 en puntos porcentuales, no niveles.\\
\var{percepcion.csv} & Seguridad percibida & Transporte y comunidad; distinto de uso de espacios.\\\bottomrule
\end{tabularx}
\end{table}

El manifiesto \var{data/infosegura_fuente.json} registra procedencia,
localizadores y versión. Los CSV conservan la precisión disponible y las
figuras no imputan años ausentes. \var{rebuild_analysis.R} regenera los contrastes.

\begin{samepage}
Reproducción desde la carpeta del proyecto, con R base:
\begin{verbatim}
Rscript --vanilla rebuild_analysis.R
\end{verbatim}
\end{samepage}
Se mantienen las series disponibles en 2021 y se separan proyecciones de
observaciones. Migración, revictimización y series mixtas se excluyen por
bases insuficientemente verificadas. No se calculan IC propios sobre estos
agregados ni se certifica comparabilidad de ponderadores entre rondas.
